\documentclass[10pt,a4paper]{article}
\usepackage[margin=0.9in]{geometry}
\usepackage[T1]{fontenc}
\usepackage[utf8]{inputenc}
\usepackage{lmodern}
\usepackage{array}
\usepackage{longtable}
\usepackage{booktabs}
\usepackage{hyperref}
\usepackage{enumitem}
\setlist{itemsep=1pt, topsep=1pt}

\title{Week 11 Project Report \\ Computer Security Software}
\author{CS 475 Introduction to Computer Security \\ Maria Nyolcas}
\date{May 2026}

\begin{document}
\maketitle

\section*{Introduction}
This report presents the Week 11 deliverable of the continuous security portfolio: a security-software taxonomy, a practical endpoint-tool evaluation, and an advisory recommendation for a university department. All activities were conducted in a legal lab setting using only safe artefacts (EICAR) and local infrastructure. Findings are mapped to the NIST Cybersecurity Framework (CSF) and CIS Controls v8.

\section*{Part A - Security Software Landscape (Theory)}

\subsection*{A-1 Major Categories and Problem Mapping}
The field of computer security software has no single universal taxonomy. The groupings below reflect both academic literature and industry practice, each addressing a distinct threat class and sitting at a specific layer of the infrastructure.

\begin{longtable}{p{3.2cm}p{3.0cm}p{2.2cm}p{5.8cm}}
\toprule
Category (acronym) & Problem solved & NIST CSF & Key limitation \\
\midrule
AV / EDR & Malicious code on endpoints & Protect, Detect, Respond & Signature lag; behavioural engines add overhead \\
Firewall (FW) & Unauthorised network access & Protect & Over-permissive rules create false assurance \\
IDS / IPS & Attack traffic detection and blocking & Detect, Respond & Rule tuning burden; encrypted traffic blind spots \\
VPN client/gateway & Eavesdropping and MITM on transit & Protect & Protects the tunnel only, not the endpoint \\
SIEM & Log aggregation and threat correlation & Identify, Detect, Respond, Recover & Dependent on log quality; alert fatigue risk \\
DLP & Data exfiltration via email / USB / cloud & Protect, Detect & Classification accuracy determines coverage \\
Vulnerability Scanner (VA) & Known vulnerabilities in software/config & Identify & Point-in-time snapshot; misses zero-days \\
PAM & Credential theft and privilege escalation & Protect, Detect & Only covers enrolled accounts; shadow admins evade \\
\bottomrule
\end{longtable}

\subsection*{A-2 Defence in Depth - Software Layer Mapping}
NIST SP 800-53 and CIS Controls v8 both prescribe layered controls rather than reliance on a single product. The table below maps security software families to their position in a defence-in-depth architecture and the rationale for each layer.

\begin{longtable}{p{2.4cm}p{5.2cm}p{6.6cm}}
\toprule
Layer & Software categories & Rationale \\
\midrule
Perimeter & NGFW, inline IPS, VPN gateway & Controls ingress/egress; forms the first boundary attackers must cross. \\
Network internal & IDS (passive tap), SIEM log ingestion, DLP proxy & Detects lateral movement and exfiltration after a perimeter bypass. \\
Host & AV/EDR, host firewall, HIDS, VA scanner & Last-resort detection and containment at the point of code execution. \\
Identity & PAM, MFA enforced by IAM tooling & Limits credential-based lateral movement and insider escalation. \\
Data & DLP agent, encryption at rest & Protects the asset even when all outer layers have been compromised. \\
\bottomrule
\end{longtable}

\noindent The core insight is that no single tool is the last line of defence. Each layer assumes the previous one has been bypassed and provides independent detection or containment capability.

\section*{Part B - Tool Evaluation and Hands-on Lab}

\subsection*{B-1 Windows Defender Evaluation Matrix}
Windows Defender (Microsoft Defender Antivirus) was evaluated against NIST SP 800-83 and CIS Control 10 criteria. It was selected because it is the default Windows endpoint protection baseline in many small organisations.

\begin{longtable}{p{4.4cm}p{1.5cm}p{8.3cm}}
\toprule
Criterion & Rating & Evidence and recommendation \\
\midrule
Signature coverage & Meets & Definitions updated multiple times daily; EICAR detected immediately. Validate update policy not delayed by GPO. \\
Behavioural / heuristic detection & Partial & Cloud heuristics (MAPS) available but require connectivity; air-gapped hosts lose this layer. Enable MAPS and sample submission. \\
Real-time (on-access) protection & Meets & EICAR blocked on file-write without manual scan trigger; Event ID 1116 logged. Enforce via Group Policy: DisableRealtimeMonitoring = 0. \\
Scheduled full scan & Partial & Default schedule exists but is configurable by users without admin friction. Enforce centrally via Intune or GPO; alert on missed scans. \\
Logging / SIEM integration & Partial & Logs to Windows Event Log (IDs 1116, 1117, 5001); WEF forwarding to SIEM requires explicit configuration. Configure WEF for Defender Operational log. \\
Tamper protection & Meets & Service stop blocked for non-admin; access denied returned on \texttt{sc stop WinDefend}. Verify via \texttt{Get-MpPreference | Select TamperProtection}. \\
Ransomware / Controlled Folder Access & Partial & Feature present but disabled by default. Enable and allow-list trusted apps before rollout to avoid workflow disruption. \\
Incident response integration & Gap & Standalone Defender lacks automated playbooks, network isolation, or ticketing without Defender for Endpoint or SOAR. Supplement with Wazuh + scripted isolation. \\
\bottomrule
\end{longtable}

\subsection*{B-2 Hands-on Detection Tests}

\textbf{Test 1 - EICAR standard test file.} The EICAR test file was written to the Desktop of a Windows 11 host with real-time protection enabled. The file was removed within approximately two seconds of creation. A Windows Security notification appeared reading \textit{``Threat found: Virus:DOS/EICAR\_Test\_File''} and Event ID 1116 was recorded in the \texttt{Microsoft-Windows-Windows Defender/Operational} event log with action \texttt{Quarantine}. Result: \textbf{PASS}.

\textbf{Test 2 - EICAR inside a ZIP archive.} The EICAR string was packaged into a \texttt{.zip} archive before writing. The archive itself was not deleted on creation. When extraction was attempted, Defender blocked the decompression and logged Event ID 1116 with the compressed file path. This confirms that Defender scans compressed content at the point of access rather than on archive creation alone. Result: \textbf{PASS}.

\textbf{Test 3 - Tamper protection (service stop attempt).} The command \texttt{sc stop WinDefend} was executed from a standard user account. The command returned \textit{``Access Denied''} (error 5) and no Event ID 5001 (service disabled) was generated. Tamper Protection prevented the service termination. Result: \textbf{PASS}.

\subsection*{B-3 Windows Firewall Rule Audit}
A review of the default Windows Firewall rule set identified the following findings relevant to the lab environment.

\begin{longtable}{p{1.5cm}p{4.6cm}p{1.6cm}p{1.4cm}p{5.1cm}}
\toprule
Dir. & Rule & State & Action & Risk commentary \\
\midrule
Inbound & Remote Desktop (TCP 3389) & Enabled & Allow & Open on all profiles; brute-force and credential-stuffing risk if internet-reachable. Scope to management source IPs. \\
Inbound & File and Printer Sharing (SMB, TCP 445) & Enabled & Allow & Acceptable on domain profile; on public profile widens lateral movement surface for ransomware. \\
Outbound & All outbound (default allow) & Enabled & Allow & Malware with C2 beaconing is unblocked. An outbound allow-list significantly limits exfiltration viability. \\
Inbound & ICMP Echo Request (ping) & Disabled & Block & Correct default on public profile; reduces host discovery exposure on untrusted networks. \\
\bottomrule
\end{longtable}

\section*{Part C - Advisory Report}
\textit{For: Head of Department, University Computing Services \\ Re: Minimum-viable security software stack for a department of approximately 50 users}

The six recommendations below are prioritised and implementable with free or open-source tools. Effort indicates expected deployment complexity.

\begin{longtable}{p{0.4cm}p{3.2cm}p{4.2cm}p{6.4cm}}
\toprule
\# & Control area & Current gap & Recommended action (tool / effort) \\
\midrule
1 & Endpoint protection & Defender enabled but unmanaged and unreported & Enforce Tamper Protection, Controlled Folder Access, and daily-update policy via Group Policy. Forward Event IDs 1116/1117 to a central log store. \textit{(Low effort)} \\
2 & Firewall hardening & RDP and outbound-all open by default & Restrict inbound RDP to management network; implement outbound allow-list for known applications. \textit{(Low effort)} \\
3 & Centralised logging & No aggregation; events isolated per host & Deploy Wazuh (free) or Elastic Security. Alert on failed logon (4625), new accounts (4720), and malware detection (1116). \textit{(Medium effort)} \\
4 & Vulnerability management & Patching occurs but no scan-based gap verification & Monthly authenticated OpenVAS scans; remediate CVSS 9+ within 72 hours, CVSS 7+ within 30 days. \textit{(Medium effort)} \\
5 & Secure remote access & Direct RDP over open internet in use & Deploy WireGuard VPN; disable external RDP; require MFA on VPN authentication. \textit{(Medium effort)} \\
6 & Privileged access & Shared local admin credentials; no session audit & Deploy Microsoft LAPS (free) for unique local admin passwords; enforce just-in-time elevation logged to SIEM. \textit{(Low effort)} \\
\bottomrule
\end{longtable}

Items 1, 2, and 6 can be implemented immediately at near-zero cost and address common breach patterns: unmonitored endpoint alerts, permissive firewall defaults, and shared credentials.

\section*{Part D - Synthesis and Reflection}
The week reinforces a recurring portfolio pattern: impactful failures are usually misconfigurations rather than exotic zero-days. Defender passed all three tests, but its value is limited without centralised logging and response. Likewise, default-allow outbound policy and open inbound RDP materially weaken perimeter security.

This week connects directly to earlier modules: Week 6 showed OSINT-driven targeting, and Week 7 showed how exposed services and weak application controls enable foothold and persistence. Security software is the operational layer that can interrupt these paths, but only when configured, monitored, and integrated.

The key conclusion is that missing governance, not missing tools, is often the real weakness. Technology without enforceable policy remains incomplete security.

\end{document}
